\documentclass[11pt,a4paper]{article}

\usepackage[utf8]{inputenc}
\usepackage[T1]{fontenc}
\usepackage[french]{babel}
\usepackage{geometry}
\geometry{margin=2.2cm}
\usepackage{hyperref}
\usepackage{graphicx}
\usepackage{array}
\usepackage{booktabs}

\title{\textbf{Rapport Synthétique -- Privacy Guard (MVP)}\\
\large Développement Assisté par IA -- Méthode Vibe Coding}

\author{Farouk \and Amazir \and Sami}
\date{Décembre 2025}

\begin{document}

\maketitle

\section{Introduction}

Privacy Guard est une application Android dont l'objectif est de protéger l'utilisateur contre les intrusions visuelles et auditives. 
L'application analyse en temps réel la caméra, l'audio, la proximité et le mouvement afin de détecter une présence potentielle, puis active automatiquement différents mécanismes de protection (blur, verrouillage, écran leurre, etc.). 

Le projet a été développé en seulement sept jours grâce à une méthodologie de développement assistée par IA, appelée \emph{Vibe Coding}, permettant une itération très rapide et une forte automatisation des étapes techniques.


\section{Fonctionnalités du MVP}

\subsection{Détection}
\begin{itemize}
    \item Camera: détection de visages (ML Kit, CameraX),
    \item Audio: détection de parole,
    \item Proximité, mouvement (accéléromètre/gyroscope),
    \item Fusion bayésienne des scores de menace.
\end{itemize}

\subsection{Protection}
Sur la base du niveau de menace, l'application active progressivement:
\begin{itemize}
    \item indicateur visuel (vert / jaune / rouge),
    \item flou progressif de l'écran,
    \item écran leurre (fake lock),
    \item verrouillage complet opaque,
    \item capture photo AES-256 en cas de tentative d'intrusion.
\end{itemize}

%\vspace{2mm}
\begin{center}
\begin{tabular}{lcc}
\toprule
\textbf{Mode} & \textbf{Seuil} & \textbf{Usage}\\
\midrule
Paranoia & 30\% & Environnements sensibles\\
Balanced & 50\% & Usage quotidien\\
Discrete & 75\% & Lieux de confiance\\
Trust Zone & 95\% & Domicile, zones sûres\\
\bottomrule
\end{tabular}
\end{center}

\section{Architecture technique}

L'architecture est structurée en quatre couches principales :

\subsection{Capteurs}
CameraSensor, AudioSensor, MotionSensor, ProximitySensor. 
Chaque capteur fournit des données traitées ensuite par un pipeline de fusion.

\subsection{Assessment}
Normalisation, scoring pondéré, fusion bayésienne, filtrage et debounce, puis calcul d'un score final de menace.

\subsection{Protection}
Affichage d'overlays système Jetpack Compose (blur, lock, fake lock) et mécanismes de capture photo chiffrée (AES-256).

\subsection{Service}
Foreground Service responsable de l'orchestration capteurs $\rightarrow$ analyse $\rightarrow$ protection (permissions, pipeline continu).




\section{Méthode Vibe Coding}

\subsection{Principes}
La méthode repose sur:
\begin{itemize}
    \item documentation exhaustive avant tout développement,
    \item prompts structurés,
    \item tests continus sur appareil réel,
    \item remontées de logs dans les prompts,
    \item commits conventionnels.
\end{itemize}

\subsection{Artefacts clés}
Dès le premier jour, plusieurs fichiers structurants ont été produits:
\texttt{SPEC.md}, \texttt{ARCHITECTURE.md}, \texttt{MVP\_ROADMAP.md}, \texttt{SECURITY\_PRIVACY.md}, etc.
Ces documents définissent l'architecture, les contraintes et les processus à respecter automatiquement lors des générations de code par l'IA.

\subsection{Résultats}
\begin{center}
\begin{tabular}{lccc}
\toprule
 & Sans IA & Vibe Coding & Gain\\
\midrule
Durée MVP & $\sim$ 4 sem. & 7 jours & 75\%\\
Bugs bloquants & $\sim$ 15 & 5 & 66\%\\
Architecture & tardive & J1 & --\\
Documentation & faible & complète & --\\
\bottomrule
\end{tabular}
\end{center}


\section{Limites observées}

\begin{itemize}
    \item calibration empirique des seuils,
    \item faux positifs audio possibles,
    \item dépendance forte aux logs (logcat),
    \item consommation énergétique élevée lorsque tous les capteurs sont actifs.
\end{itemize}


\section{Évolutions prévues}

Six spécifications post-MVP ont été planifiées:
\begin{itemize}
    \item v1.5 : zones de confiance (GPS, WiFi),
    \item v2.0 : apprentissage comportemental,
    \item v2.5 : mode furtif et panic mode,
    \item v3.0 : version entreprise (MDM).
\end{itemize}


\section{Conclusion}

Le projet Privacy Guard démontre que le développement assisté par IA peut réduire drastiquement le temps de réalisation d'un MVP complexe.
La qualité de l'architecture, la rapidité d'itération et la documentation permanente ont permis de livrer en une semaine un produit fonctionnel et extensible.
Cette approche positionne l'IA comme un véritable \emph{co-développeur senior} lorsque le contexte et la documentation sont correctement fournis.

\vspace{4mm}
\noindent\textit{Projet développé en novembre--décembre 2025 en Vibe Coding avec Claude Sonnet 4.5.}

\end{document}

